\documentclass[11pt]{article}

\usepackage[margin=1in]{geometry}
\usepackage{amsmath,amssymb,amsthm}
\usepackage{algorithm}
\usepackage{algpseudocode}
\usepackage{enumitem}
\usepackage{hyperref}
\usepackage{xcolor}
\usepackage{tikz}
\usepackage{booktabs}

\hypersetup{colorlinks=true, linkcolor=blue, urlcolor=blue, citecolor=blue}

\newtheorem{theorem}{Theorem}[section]
\newtheorem{lemma}[theorem]{Lemma}
\newtheorem{proposition}[theorem]{Proposition}
\newtheorem{corollary}[theorem]{Corollary}
\newtheorem{definition}[theorem]{Definition}
\newtheorem{example}[theorem]{Example}

\title{TOAE209/TOAH209 -- Design and Analysis of Algorithms\\[4pt]
\large Week 8: Dynamic Programming I: Weighted Interval Scheduling and Knapsack}
\author{Foreign Trade University}
\date{}

\begin{document}
\maketitle
\tableofcontents

\section{Introduction: The Dynamic Programming Paradigm}

This week we begin our study of \textbf{dynamic programming} (DP), arguably the most
powerful general-purpose algorithm-design technique in this course. Where greedy
algorithms (Weeks 4--5) commit to a single locally-optimal choice at each step and never
look back, dynamic programming instead \emph{considers all the choices available at each
step}, but does so efficiently by reusing the answers to subproblems it has already
solved.

\medskip
\noindent\textit{(Note: this is the Midterm week -- the material below also serves as a
consolidation of the problem-solving habits developed so far: define the object precisely,
write a recurrence, prove it correct, analyze its cost, and verify empirically.)}

The name ``dynamic programming'' is historical (Richard Bellman, 1950s; ``programming''
here means ``planning,'' not writing code). The essential idea is captured by one slogan:

\begin{quote}
\emph{Solve each subproblem once, store its answer, and build the answer to a larger
problem from the answers to smaller ones.}
\end{quote}

\subsection{When does dynamic programming apply?}

Two structural properties must hold for a problem to admit an efficient DP solution.

\begin{definition}[Optimal substructure]
A problem has \emph{optimal substructure} if an optimal solution to the whole problem can
be expressed in terms of optimal solutions to its subproblems. Equivalently: once we fix
one ``decision'' (e.g.\ whether to include the last item), what remains is an independent,
smaller instance of the same problem, and the overall optimum uses an optimum of that
smaller instance.
\end{definition}

\begin{definition}[Overlapping subproblems]
A problem has \emph{overlapping subproblems} if a naive recursive solution would solve the
\emph{same} subproblem many times. This is what distinguishes dynamic programming from
plain divide-and-conquer (where subproblems are typically disjoint, as in merge sort).
\end{definition}

When both hold, we can write a \textbf{recurrence} (the Bellman equation) expressing the
optimum of an instance in terms of the optima of smaller instances, and then evaluate it
without recomputation. The textbook example of overlapping subproblems is the Fibonacci
recurrence $F(n) = F(n-1) + F(n-2)$: the naive recursion makes $2F(n+1)-1$ calls
(exponentially many), whereas remembering each $F(k)$ reduces this to $n+1$ distinct
subproblems. Problem 12 asks you to instrument exactly this blow-up and confirm that DP
does linear work.

\subsection{Two implementation styles: memoization vs.\ bottom-up}

There are two equivalent ways to evaluate a DP recurrence without recomputation.

\begin{itemize}
    \item \textbf{Top-down (memoization).} Write the recurrence as a recursive function,
    but cache (``memoize'') the result of each subproblem the first time it is computed; on
    subsequent calls, return the cached value immediately. Subproblems are solved
    \emph{lazily}: only those actually reached are ever computed. This is often the easiest
    to write (it mirrors the recurrence directly) and avoids computing unreachable
    subproblems.

    \item \textbf{Bottom-up (tabulation).} Identify an order in which subproblems depend
    only on \emph{smaller} ones, then fill a table iteratively in that order, so that
    whenever we need a subproblem's value it is already in the table. This avoids recursion
    overhead, makes the running time and space transparent, and often enables space
    optimizations (e.g.\ keeping only the last row of a 2-D table).
\end{itemize}

Both compute the same values and have the same asymptotic running time (number of distinct
subproblems times the work per subproblem). Problem 4 asks you to implement a DP both ways
and confirm they agree. We will generally present the bottom-up form for analysis and the
top-down form when the recurrence is more natural recursively.

\subsection{A recipe for designing a DP}

For every DP this week (and in Week 9), we follow the same five steps:
\begin{enumerate}[label=(\arabic*)]
    \item \textbf{Define the subproblem} precisely, in words and as an array
    (e.g.\ ``$M[j]$ = optimal value over the first $j$ intervals'').
    \item \textbf{Write the recurrence} relating a subproblem to smaller ones, including
    the \textbf{decision} being made (take vs.\ skip, cut here vs.\ later).
    \item \textbf{State the base cases.}
    \item \textbf{Choose an evaluation order} (bottom-up) or memoize (top-down), and
    analyze the running time as (number of subproblems) $\times$ (work per subproblem).
    \item \textbf{Reconstruct an optimal solution} by tracing back through the table, not
    just reporting the optimal \emph{value}.
\end{enumerate}

\section{Weighted Interval Scheduling}

\subsection{Problem statement}

Recall the (unweighted) Interval Scheduling problem from Week 4, where we maximized the
\emph{number} of compatible intervals and the earliest-finish-time greedy algorithm was
optimal. We now attach a \textbf{value} to each interval.

\begin{definition}[Weighted Interval Scheduling]
We are given $n$ requests, where request $i$ has a start time $s_i$, a finish time $f_i$
(with $s_i < f_i$), and a \emph{weight} (or value) $w_i \ge 0$. Two requests are
\emph{compatible} if their half-open intervals $[s_i, f_i)$ and $[s_j, f_j)$ do not
overlap. The goal is to select a set of mutually compatible requests of \textbf{maximum
total weight}.
\end{definition}

\subsection{Why greedy fails, and what replaces it}

The earliest-finish-time greedy rule maximizes the \emph{count} of intervals, but a single
high-value interval can be worth more than many low-value ones it conflicts with. Sorting
by weight fails too: a very heavy interval might block two even-heavier compatible ones. No
simple greedy ordering is correct -- so we turn to dynamic programming.

Sort the requests by finish time, so $f_1 \le f_2 \le \cdots \le f_n$. The key quantity is:

\begin{definition}[The function $p(j)$]
For each request $j$, let $p(j)$ be the largest index $i < j$ such that request $i$ is
compatible with request $j$ (i.e.\ $f_i \le s_j$), or $0$ (``none'') if no such $i$
exists.
\end{definition}

Because the requests are sorted by finish time, $p(j)$ can be found by \textbf{binary
search} on the array of finish times: $p(j)$ is the rightmost interval whose finish time is
$\le s_j$. Computing all of $p(\cdot)$ thus takes $O(n \log n)$. Problem 1 asks you to
implement this.

\subsection{The recurrence}

Consider the last interval, request $n$ (the one finishing last). Any optimal solution
either \emph{includes} it or not -- and this single binary decision is the heart of the
DP.

\begin{itemize}
    \item If the optimum \textbf{includes} $n$, then it can contain no interval that
    overlaps $n$; the remaining intervals it may use are exactly $1, \ldots, p(n)$. So the
    optimum is $w_n$ plus the optimum over $\{1, \ldots, p(n)\}$.
    \item If the optimum \textbf{excludes} $n$, then it is simply the optimum over
    $\{1, \ldots, n-1\}$.
\end{itemize}

Let $M[j]$ denote the maximum total weight obtainable using only requests $1, \ldots, j$
(after sorting by finish time). Then:
\[
M[0] = 0, \qquad
M[j] = \max\bigl(\, M[j-1], \;\; w_j + M[p(j)] \,\bigr).
\]

\begin{theorem}[Correctness of the recurrence]
For every $j$, $M[j]$ as defined above equals the maximum total weight of a compatible
subset of $\{1, \ldots, j\}$.
\end{theorem}

\begin{proof}
By strong induction on $j$. The base case $M[0]=0$ is the empty instance. For $j \ge 1$,
let $O_j$ be an optimal compatible subset of $\{1,\ldots,j\}$. Either $j \in O_j$ or not.
If $j \notin O_j$, then $O_j$ is a compatible subset of $\{1,\ldots,j-1\}$, so its weight
is at most $M[j-1]$ (by the inductive hypothesis), and conversely any compatible subset of
$\{1,\ldots,j-1\}$ is also valid for $\{1,\ldots,j\}$, so this branch contributes exactly
$M[j-1]$. If $j \in O_j$, then no interval in $O_j$ overlaps $j$, so every other member has
index $\le p(j)$; thus $O_j \setminus \{j\}$ is a compatible subset of $\{1,\ldots,p(j)\}$
of weight $M[j] - w_j$, which by induction is at most $M[p(j)]$, giving weight at most
$w_j + M[p(j)]$; and this value is achievable (take an optimum of $\{1,\ldots,p(j)\}$ plus
$j$). The optimum is the better of the two branches, which is exactly the stated maximum.
\end{proof}

\subsection{Algorithm and running time}

\begin{algorithm}[h]
\caption{Weighted-Interval-Scheduling (bottom-up)}
\begin{algorithmic}[1]
\State Sort requests by finish time so that $f_1 \le \cdots \le f_n$
\State Compute $p(1), \ldots, p(n)$ by binary search on finish times
\State $M[0] \gets 0$
\For{$j = 1$ to $n$}
    \State $M[j] \gets \max\bigl(M[j-1],\; w_j + M[p(j)]\bigr)$
\EndFor
\State \Return $M[n]$
\end{algorithmic}
\end{algorithm}

Sorting is $O(n\log n)$; computing all $p(j)$ is $O(n \log n)$; the main loop is $O(n)$.
Total: $O(n \log n)$. Problems 2 (value) and 4 (top-down memoization) implement this.

\subsection{Reconstructing an optimal solution}

The array $M$ gives the optimal \emph{value}; to recover the optimal \emph{set} we trace
back. Starting at $j = n$, request $j$ is in the optimum iff the ``take'' branch was at
least as good as the ``skip'' branch:
\[
w_j + M[p(j)] \;\ge\; M[j-1].
\]
If so, output $j$ and jump to $j \gets p(j)$; otherwise move to $j \gets j-1$. This runs in
$O(n)$. Problem 3 asks you to implement this traceback and verify that the reconstructed
set is feasible and has weight exactly $M[n]$.

\section{Subset Sum and the 0/1 Knapsack Problem}

\subsection{Subset Sum}

\begin{definition}[Subset Sum]
Given $n$ non-negative integers $a_1, \ldots, a_n$ and a target $T$, decide whether some
subset of the $a_i$ sums to exactly $T$.
\end{definition}

Let $R[i][t]$ be true iff some subset of the first $i$ numbers sums to $t$. Considering
item $i$ (take it or skip it):
\[
R[0][0] = \text{true}, \quad R[0][t] = \text{false } (t>0), \qquad
R[i][t] = R[i-1][t] \;\lor\; \bigl(t \ge a_i \;\land\; R[i-1][t - a_i]\bigr).
\]
The answer is $R[n][T]$. The table has $(n+1)(T+1)$ Boolean entries, each filled in $O(1)$,
so the running time is $O(nT)$. Problem 5 asks you to implement this, using the standard
1-D optimization: keep a single Boolean array \texttt{reachable} of length $T+1$ and, for
each $a_i$, update it for $t$ from $T$ \emph{down to} $a_i$. The descending order is
essential -- it ensures each item is used \emph{at most once} (if we updated in ascending
order, $a_i$ could be reused, which would instead solve the \emph{unbounded} version).

\subsection{0/1 Knapsack}

\begin{definition}[0/1 Knapsack]
Given $n$ items, where item $i$ has an integer weight $w_i \ge 0$ and value $v_i \ge 0$,
and a knapsack of integer capacity $W$, select a subset of items of \emph{maximum total
value} whose total weight is at most $W$. Each item may be taken at most once (hence
``0/1'').
\end{definition}

Subset Sum is the special case $v_i = w_i$, $W = T$, asking only for feasibility. The DP is
the natural generalization. Let $\mathit{dp}[i][w]$ be the maximum value achievable using
the first $i$ items with capacity $w$:
\[
\mathit{dp}[0][w] = 0, \qquad
\mathit{dp}[i][w] =
\begin{cases}
\mathit{dp}[i-1][w] & \text{if } w_i > w,\\[2pt]
\max\bigl(\mathit{dp}[i-1][w],\; v_i + \mathit{dp}[i-1][w - w_i]\bigr) & \text{otherwise.}
\end{cases}
\]
The first case says: if item $i$ does not fit, we must skip it. The second case is the
take-vs-skip decision. The answer is $\mathit{dp}[n][W]$.

\begin{algorithm}[h]
\caption{Knapsack-0/1 (bottom-up, 2-D table)}
\begin{algorithmic}[1]
\For{$w = 0$ to $W$} \State $\mathit{dp}[0][w] \gets 0$ \EndFor
\For{$i = 1$ to $n$}
    \For{$w = 0$ to $W$}
        \State $\mathit{dp}[i][w] \gets \mathit{dp}[i-1][w]$
        \If{$w_i \le w$}
            \State $\mathit{dp}[i][w] \gets \max\bigl(\mathit{dp}[i][w],\; v_i + \mathit{dp}[i-1][w - w_i]\bigr)$
        \EndIf
    \EndFor
\EndFor
\State \Return $\mathit{dp}[n][W]$
\end{algorithmic}
\end{algorithm}

\begin{theorem}
Algorithm Knapsack-0/1 computes the optimal value in $O(nW)$ time and $O(nW)$ space.
\end{theorem}

\begin{proof}
Correctness follows from the recurrence by induction on $i$, exactly as for Weighted
Interval Scheduling: $\mathit{dp}[i][w]$ considers the binary decision on item $i$ and
defers the rest to optimal sub-instances $\mathit{dp}[i-1][\cdot]$ (optimal substructure).
The table has $(n+1)(W+1)$ entries, each computed in $O(1)$, giving $O(nW)$.
\end{proof}

Problem 6 implements the value; Problem 7 reconstructs the chosen items by tracing back:
item $i$ was taken iff $\mathit{dp}[i][w] \ne \mathit{dp}[i-1][w]$, in which case subtract
$w_i$ from $w$ and continue to row $i-1$. (A standard space optimization keeps only one
row, but then reconstruction requires either storing decisions or a divide-and-conquer
trick; we keep the full table for clarity.)

\subsection{Why $O(nW)$ is only \emph{pseudo}-polynomial}

The running time $O(nW)$ looks polynomial, and indeed it is polynomial in the
\emph{numeric value} $W$. But the \emph{size of the input} -- the number of bits needed to
write it down -- is only $O(n \log W)$, because the capacity $W$ is written in binary using
about $\log_2 W$ bits. So $W$ itself is \emph{exponential} in the input size: doubling the
number of bits of $W$ squares $W$, and hence squares the running time. An algorithm whose
running time is polynomial in the numeric values of the inputs, but exponential in their
bit-length, is called \textbf{pseudo-polynomial}.

This matters because 0/1 Knapsack is \textbf{NP-hard}: no genuinely polynomial-time (in
the input size) algorithm is known, and one is believed not to exist. The $O(nW)$ DP is
efficient \emph{only when $W$ is small} (e.g.\ bounded by a polynomial in $n$). For
instances with huge capacities, the table is astronomically large. This is our first
concrete encounter with the boundary between tractable and intractable problems, a theme
we will return to in the complexity unit.

\section{Unbounded Knapsack: Coin Change and Rod Cutting}

In the \emph{unbounded} knapsack variant, each item type may be used \emph{any number of
times}. Two classic instances appear in this week's exercises.

\subsection{Coin change: minimum number of coins}

\begin{definition}[Coin Change -- minimization]
Given coin denominations $c_1, \ldots, c_k$ (unlimited supply of each) and an amount $A$,
find the minimum number of coins summing to exactly $A$, or report that $A$ cannot be
made.
\end{definition}

Let $\mathit{dp}[a]$ be the minimum number of coins summing to $a$:
\[
\mathit{dp}[0] = 0, \qquad
\mathit{dp}[a] = 1 + \min_{\,c_i \le a} \mathit{dp}[a - c_i]
\quad(\text{or } +\infty \text{ if no } c_i \le a \text{ helps}).
\]
The answer is $\mathit{dp}[A]$, or $-1$ if it is $+\infty$. The table has $A+1$ entries,
each requiring an $O(k)$ minimization, so the running time is $O(kA)$ -- again
pseudo-polynomial. Note the crucial difference from Subset Sum: because coins are
\emph{reusable}, we iterate $a$ in \emph{ascending} order and do \emph{not} reverse the
loop. Unlike the Week 4 greedy coin-change algorithm (correct only for ``canonical'' coin
systems), this DP is correct for \emph{any} coin system. Problem 8 implements it.

\subsection{Coin change: number of ways}

\begin{definition}[Coin Change -- counting]
With the same input, count the number of distinct \emph{combinations} (multisets, where
order does not matter) of coins summing to $A$.
\end{definition}

Let $\mathit{ways}[a]$ be the number of combinations summing to $a$, with $\mathit{ways}[0]
= 1$ (the empty combination). To count \emph{combinations} rather than \emph{ordered
sequences}, we loop over the coins on the \emph{outside} and amounts on the inside:
\[
\textbf{for } c \text{ in coins:}\quad \textbf{for } a = c \text{ to } A:\quad
\mathit{ways}[a] \mathrel{+}= \mathit{ways}[a - c].
\]
Putting the coin loop outside fixes a canonical order in which coins are ``committed,'' so
each combination is counted exactly once. Swapping the loop order would instead count
permutations. Problem 9 implements this and is the same idea behind LeetCode's
\emph{Coin Change II}.

\subsection{Rod cutting}

\begin{definition}[Rod Cutting]
Given a rod of integer length $n$ and a price table where $p_\ell$ is the price of a piece
of length $\ell$, cut the rod into integer pieces to maximize total revenue.
\end{definition}

This is unbounded knapsack with ``item $\ell$'' = a piece of length $\ell$, weight $\ell$,
value $p_\ell$, and capacity $n$ (each length usable repeatedly). Let $\mathit{rev}[L]$ be
the maximum revenue from a rod of length $L$:
\[
\mathit{rev}[0] = 0, \qquad
\mathit{rev}[L] = \max_{1 \le \ell \le L}\bigl(p_\ell + \mathit{rev}[L - \ell]\bigr).
\]
The running time is $O(n^2)$. Recording, for each $L$, the cut length $\ell$ achieving the
maximum lets us reconstruct the list of pieces. Problem 10 implements both the revenue and
the cut list.

\section{Longest Increasing Subsequence}

As a final example of a one-dimensional DP with a slightly different flavor:

\begin{definition}[Longest Increasing Subsequence (LIS)]
Given a sequence $a_1, \ldots, a_n$, find the length of the longest subsequence (not
necessarily contiguous, but keeping the original order) that is strictly increasing.
\end{definition}

Let $L[i]$ be the length of the longest increasing subsequence \emph{ending at} index $i$:
\[
L[i] = 1 + \max\bigl(\{\,L[j] : j < i,\; a_j < a_i\,\} \cup \{0\}\bigr).
\]
The answer is $\max_i L[i]$ (or $0$ if the sequence is empty). Each $L[i]$ scans all
earlier indices, giving $O(n^2)$. (An $O(n \log n)$ patience-sorting algorithm also exists;
we defer it.) Problem 11 implements the $O(n^2)$ DP and verifies it against a brute-force
search over all subsequences on small inputs -- a good template for empirically validating
any DP.

\section{Summary and Looking Ahead}

This week introduced \textbf{dynamic programming} as a systematic technique for problems
with \emph{optimal substructure} and \emph{overlapping subproblems}, contrasting
\emph{top-down memoization} with \emph{bottom-up tabulation}. The case studies were:

\begin{itemize}
    \item \textbf{Weighted Interval Scheduling}: the recurrence $M[j] = \max(M[j-1],\,
    w_j + M[p(j)])$, with $p(j)$ found by binary search, solved in $O(n \log n)$, plus
    solution reconstruction by traceback.
    \item \textbf{Subset Sum and 0/1 Knapsack}: the 2-D table $\mathit{dp}[i][w]$ in
    $O(nW)$ time, with item reconstruction -- and the observation that $O(nW)$ is only
    \emph{pseudo-polynomial}, since $W$ is exponential in the input's bit-length (Knapsack
    is NP-hard).
    \item \textbf{Unbounded knapsack}: coin change (minimum coins, and counting
    combinations) and rod cutting, where items may be reused -- distinguished from 0/1 by
    the loop direction/order.
    \item \textbf{Longest Increasing Subsequence}: an $O(n^2)$ DP, verified against brute
    force.
\end{itemize}

The recurring discipline is the five-step recipe of Section 1.4: define the subproblem,
write the recurrence (centered on a single decision), state base cases, evaluate
efficiently, and reconstruct the optimal solution -- not just its value.

Next week (Dynamic Programming II, Kleinberg \& Tardos Ch.\ 6 continued) we extend this
toolkit to \emph{two-dimensional} alignment problems -- \textbf{sequence alignment / edit
distance}, the basis of spell-checkers and computational biology -- and to
\textbf{shortest-path} problems on graphs with negative edges via the Bellman--Ford and
Floyd--Warshall dynamic programs, where the subproblems are indexed by ``number of edges
allowed'' or ``set of intermediate vertices.''

\end{document}
